%% kapa: exam=1 round=19 year=2008 period=1 form=A subject=civil_law subject_ko=민법 seq=11 number=11
%% source: Q-Net 감정평가사 기출문제 (cst00309) artlSeq=1001962 file=extracted/1008/1001962/2008년 제19회 감정평가사 1차시험 기출문제/1교시.hwp
%% conversion: kordoc 4.24.51 -> tools/gen1.py (auto); review=proofread-sonnet; flags=none
\begin{question}{11}
물건에 관한 설명으로 틀린 것은? (다툼이 있는 경우에는 판례에 의함)
\begin{choices}
  \item 장소․종류․수량지정 등의 방법으로 특정할 수 있다면 증감변동하는 동산도 그 전부를 하나의 물건으로 보아 담보의 대상으로 할 수 있다.
  \item 건축주가 사회통념상 독립한 건물이라고 볼 수 있는 형태와 구조를 갖추고 있는 미완성 건물을 타인에게 매도한 후 건축주 명의변경절차를 마쳤더라도 원래의 건축주가 건물을 원시취득한다.
  \item 온천수는 토지의 구성부분이다.
  \item 건물의 개수(個數)는 거래관념에 따라 정해지는 것이 아니라, 최소한의 기둥과 지붕 그리고 주벽의 설치 여부와 같은 물리적 구조에 의하여 결정된다.
  \item 토지소유자의 승낙을 얻지 않고 권원 없이 그 토지 위에 식재한 수목은 토지에 부합한다.
\end{choices}
\end{question}
\answer{4}
